
\chapter{Technical Analysis}

\section{Member representation}

In this chapter, we will outline several ways how to represent type members in the code.

The first idea is to use the objects provided by Reflection, such as \textit{MethodInfo}, which provide the context data we need.
As mentioned before, these objects contain properties such as \texttt{IsStatic}, \texttt{IsPublic}, etc.

However, as mentioned before, the documentation comments are not part of the assembly. 
Therefore, the objects provided by reflection do not contain any data regarding the documentation comments.

Option 1: reflection provided objects + external XML file

Option 2: tuple (object, doc comment)

disadvantages: 

Option 3: custom objects (i.e. MethodData)

advantages: common interfaces, doc comments included

Nice object design, great for extensibility -> the users extending the library may easily 

\begin{itemize}
    \item Data objects vs custom classes
\end{itemize}
Advantages:
\begin{itemize}
    \item Common interfaces
    \item Documentation comments included
\end{itemize}

describe MemberID

\section{Doc comment extraction}

- how to add comments to the members

\section{InheritDoc resolving}

- how to resolve inheritdoc
- three types of inheritance: type, member and code reference

\begin{itemize}
    \item Depth-first search graph
\end{itemize}

\section{Doc Comment Parser}
\begin{itemize}
    \item Describe XML documentation tags converted to HTML.
\end{itemize}

\section{Template generator}
\begin{itemize}
    \item Describe XML documentation tags converted to HTML.
\end{itemize}

\section{Template Language Choice}

\subsection{No template language?}

\subsection{Mustache}
\begin{itemize}
    \item Stubble.NET library
\end{itemize}

\subsection{Handlebars}
\begin{itemize}
    \item Handlebars.NET library
\end{itemize}

\subsection{Liquid}
\begin{itemize}
    \item DotLiquid
    \item Fluid.NET
    \item Scriban
\end{itemize}

\subsection{Razor}
\begin{itemize}
    \item RazorLight
    \item Razor components in Console apps
\end{itemize}

\subsection{Conclusion}

\section{Templates}

\subsection{Template models}
- describe
- template model x memberData objects

\subsection{Template model creators}
- the conversion memberData -> template model

\subsection{Razor templates}
- decomposition
